\subsection*{Centerline precision and soundness of a pulse certificate}
\label{subsec:centerline-proof-gate}
The source's normalized pulse denominator is
\[
D=\tfrac12-h,\quad\beta=\tfrac12-\lambda,\quad
c(\eta)=\frac{2D\eta^2}{1+\eta^2},\quad
q_0(\eta)=\frac{\lambda-h+c(\eta)}{1-\lambda}.
\]
Thus, for $0<h\le\lambda/10^4$ and $0<\lambda\le10^{-14}$,
\[
\frac{9999}{10000}\lambda\le q_0(0)
=\frac{\lambda-h}{1-\lambda}
\le\frac{\lambda}{1-\lambda},\qquad q_0(\eta)\ge q_0(0)>0.
\]
\begin{proposition}[Necessary absolute source precision]
\label{prop:source-pulse-precision}
If the A.29--A.30 comparison inequalities
$|e_Q/q_0|\le10^{-10}$ and $|e_N/q_0|\le10^{-7}$
hold, then
\[
|e_Q(0)|\le 10^{-10}\frac{\lambda-h}{1-\lambda},
\qquad |e_N(0)|\le10^{-7}\frac{\lambda-h}{1-\lambda}.
\]
At $\lambda=10^{-14}$ and $h=\lambda/10^4$, the
right sides are of order $10^{-24}$ and $10^{-21}$,
respectively. They shrink proportionally for smaller $\lambda$.
\end{proposition}
\begin{proof}
Use $c\ge0$ with equality at $\eta=0$; multiply the
relative inequalities by the exact positive $q_0(0)$.
The endpoint estimates follow by rational arithmetic.
\end{proof}
Thus an absolute error of merely $10^{-16}$ obtained by
subtracting order-one source terms is inadequate. A correct
interval computation must first exploit source cancellations
and use sufficiently accurate outward enclosures.

\paragraph{Cancellation before interval evaluation.}
Set
\[
S=D+c,\quad
C_d=\frac{\lambda S(1-\lambda)}
{\beta^2(\lambda-h+c)},\quad w_0=2R-C_dR',
\]
\[
e_m=m-\frac R\beta+\frac{\lambda R'}{\beta^2},
\quad e_Q=Q_s-q_0,\quad
e_N=\frac{N_s}{E}+R-Sm+D\eta m_\eta.
\]
\begin{lemma}[Exact denominator-cleared pulse identity]
\label{lem:pulse-source-cleared}
Whenever the physical denominators are nonzero,
\begin{align}
\beta^2(N_s/E-w_0Q_s)={}&
q_0\lambda R\beta+S\beta^2 e_m
-D\eta\beta^2m_\eta+\beta^2e_N-\beta^2w_0e_Q.
\label{eq:source-cleared-pulse}
\end{align}
Equivalently,
\[
\frac{N_s}{EQ_s}-w_0=
\frac{q_0(\lambda/\beta)R+Se_m-D\eta m_\eta+
e_N-w_0e_Q}{q_0+e_Q}.
\]
\end{lemma}
\begin{proof}
After multiplying by $\beta^2$, the coefficient of $R$
is $\beta[S-\beta-q_0(2\beta+\lambda)]$; it vanishes since
$S-\beta=q_0(1-\lambda)$ and $2\beta+2\lambda=1$.
The $R'$ coefficient is
$C_dq_0\beta^2-\lambda S=0$. The remaining
terms cancel by definition. Independent symbolic and
rational replays test both identities.
\end{proof}

\paragraph{Authenticating proof inputs rather than labels.}
The existing rational full-cover program checks whether the
provided rectangles cover $[0,13]\times[-1,1]$ and whether
their supplied intervals meet each residual budget. Its
former manifest accepted any syntactically valid SHA-256
string without comparing it to file bytes. The new
\texttt{pulse\_provenance\_replay\_gate.py} binds the three
primary evidence files to recomputed SHA-256 digests, rejects
missing or tampered files and traversal outside the evidence
root, and then runs the exact rectangle and residual checks.
Negative tests also reject empty coverings.

This improves reproducibility but does not prove that the
submitted intervals enclose the actual Appendix-A fields.
Such a theorem must independently validate, on every closed
cell, the source's field construction, derivatives, ODE
comparison and interval arithmetic. No complete source-derived
interval dataset or jointly admissible source tuple is present
in this release. A matching hash is never a substitute for
mathematical enclosure validity.

\paragraph{Lean and physical status.}
\texttt{PulseResidualPrecision.lean} proposes kernel lemmas for
the cleared polynomial identity and centerline rational
scales. Until a Lean job completes successfully, those
statements are not kernel-certified. The full source interval
construction, exterior stress cone, nonlinear Appendix-B
analytic profile, unforced correction hierarchy and exact
physical zero-force Navier--Stokes residual remain unproved.
The originating OpenAI whole-space theorem uses a smooth
external force \cite{openai2026}; it does not prove the
proposed unforced Theorem 1.1.
